\section{Conclusion}
In this work, we conducted a comprehensive empirical study to characterize the scaling behaviors of reinforcement learning post-training for large language models. Our analysis yields three fundamental findings: 

First, we establish a predictive power-law formulation that models the log-linear relationship between test loss and resource consumption. We demonstrate that this formulation enables reliable performance forecasting, accurately predicting the learning efficiency of larger model and estimating final trajectories from early training dynamics. 

Second, we identify a inherent saturation trend in learning efficiency. While larger models consistently exhibit superior compute and sample efficiency, the marginal gains in the learning efficiency coefficient $k(N)$ do not grow indefinitely; instead, they diminish as model size increases, asymptotically approaching a theoretical limit. 

Third, we find that RL model performance in data-constrained regimes is primarily governed by the total volume of training data rather than sample uniqueness. Our experiments validate that moderate data reuse is a highly effective strategy, where increasing the repetition factor allows models to maintain performance improvements without requiring additional unique samples.